% Part: lambda-calculus
% Chapter: church-rosser
% Section: beta-reduction

\documentclass[../../../include/open-logic-section]{subfiles}

\begin{document}

\olfileid{lam}{cr}{b}

\olsection{$\beta$-न्यूनीकरण}

\begin{lem}\ollabel{lem:one-par}
  $M \bredone M'$ असेल, तर $M \bredpar M'$.
\end{lem}
\begin{proof} $M \bredone M'$ च्या निष्पत्तीवर विगमन करू.
  मूळस्थानी संकुचनाच्या प्रकरणात काही $x$, $N$ आणि~$Q$ साठी
  $M$ हे $(\lambd[x][N])Q$ आणि $M'$ हे
  $\Subst{N}{Q}{x}$ आहे. \olref[pb]{thm:refl} नुसार
  $N \bredpar N$ आणि $Q \bredpar Q$ असल्याने
  \olref[pb]{defn:bredpar}\olref[pb]{defn:bredpar4} नुसार
  $(\lambd[x][N])Q \bredpar \Subst{N}{Q}{x}$ लगेच मिळते.
  अमूर्तीकरण किंवा उपयोजनाच्या आत संकुचन झाल्यास,
  विगमन गृहीतक, अपरिवर्तित उपपदाची समांतर परावर्तिता
  आणि संबंधित रचना-नियमांवरून निष्कर्ष मिळतो.
\end{proof}

\begin{lem}\ollabel{lem:par-red}
  $M \bredpar M'$ असेल, तर $M \bred M'$.
\end{lem}

\begin{proof} $M \bredpar M'$ च्या !!{derivation} वर विगमन करू.
  \begin{enumerate}
  \item अखेरचा नियम \olref[pb]{defn:bredpar1} आहे: मग $M$
    आणि $M'$ दोन्ही $x$ आहेत आणि $x \bred x$.
  \item अखेरचा नियम \olref[pb]{defn:bredpar2} आहे: काही
    $x$, $N$, $N'$ साठी $M$ हे $\lambd[x][N]$
    आणि $M'$ हे $\lambd[x][N']$ आहे, जेथे
    $N \bredpar N'$. विगमन गृहीतकानुसार $N \bred N'$.
    मग $N \bred N'$ साठी लागणारी तीच $\bredone$
    संकुचने अमूर्तीकरणाच्या आत करून
    $\lambd[x][N] \bred \lambd[x][N']$ मिळते.
  \item अखेरचा नियम \olref[pb]{defn:bredpar3} आहे: काही
    $P$, $Q$, $P'$, $Q'$ साठी $M$ हे $PQ$ आणि
    $M'$ हे $P'Q'$ आहे, जेथे $P \bredpar P'$
    आणि $Q \bredpar Q'$. विगमन गृहीतकानुसार
    $P \bred P'$ आणि $Q \bred Q'$. म्हणून आधी
    $P \bred P'$ आणि नंतर $Q \bred Q'$ असे
    न्यूनीकरण केल्यावर $PQ \bred P'Q'$ मिळते.
  \item अखेरचा नियम \olref[pb]{defn:bredpar4} आहे: काही
    $x$, $N$, $N'$, $Q$, $Q'$ साठी $M$ हे
    $(\lambd[x][N])Q$ आणि $M'$ हे
    $\Subst{N'}{Q'}{x}$ आहे, जेथे $N \bredpar N'$
    आणि $Q \bredpar Q'$. विगमन गृहीतकानुसार
    $Q \bred Q'$ आणि $N \bred N'$. म्हणून आधी
    $N \bred N'$, मग $Q \bred Q'$, आणि शेवटी
    $(\lambd[x][N'])Q'$ चे $\Subst{N'}{Q'}{x}$
    मध्ये संकुचन करून
    $(\lambd[x][N])Q \bred \Subst{N'}{Q'}{x}$ मिळते.
  \end{enumerate}
\end{proof}

\begin{lem}\ollabel{lem:str}
  $\bredpar$ समाविष्ट करणारा सर्वांत लहान संक्रमक
  संबंध $\bred$ आहे.
\end{lem}

\begin{proof}
  $\bredpar$ समाविष्ट करणारा सर्वांत लहान संक्रमक
  संबंध $\xred$ मानू.

  $\bred \subseteq \xred$: समजा $M \bred M'$,
  म्हणजे $M \ident M_1 \bredone \dots \bredone M_k \ident M'$.
  \olref{lem:one-par} नुसार $M \ident M_1 \bredpar
  \dots \bredpar M_k \ident M'$. $\xred$ मध्ये
  $\bredpar$ समाविष्ट आहे आणि तो संक्रमक आहे;
  म्हणून $M \xred M'$. शून्य संकुचनांच्या प्रकरणातही
  समांतर न्यूनीकरणाच्या परावर्ती गुणधर्मामुळे हाच निष्कर्ष मिळतो.

  $\xred \subseteq \bred$: समजा $M \xred M'$,
  म्हणजे $M \ident M_1 \bredpar \dots \bredpar M_k \ident M'$.
  \olref{lem:par-red} नुसार $M \ident M_1 \bred
  \dots \bred M_k \ident M'$. $\bred$ संक्रमक
  असल्याने $M \bred M'$.
\end{proof}

\begin{thm}\ollabel{thm:cr}
  $\bred$ ला चर्च–रॉसर गुणधर्म आहे.
\end{thm}

\begin{proof}
  \olref[dap]{thm:str}, \olref[pb]{thm:cr} आणि
  \olref{lem:str} वरून लगेच मिळते.
\end{proof}

\end{document}
